\appendix
\section{The DT-margin Heuristic Baseline}
\label{app:dtmargin}
%==============================================================================

VA-ATACOM's velocity augmentation has a heuristic predecessor we call the \textbf{DT-margin} filter --- the ``DT-margin (heuristic)'' row of Table~\ref{tab:main}. It inflates the keepout radius linearly with speed, $r_\text{eff}(\boldsymbol{v}) = r_\text{base}(1+\alpha\|\boldsymbol{v}\|)$, projects out the inward radial component, and adds the same sim-BRT reward shaping. As shown in \S\ref{sec:design}, this linear margin is the \emph{first-order} approximation of VA-ATACOM's quadratic braking constraint (\ref{eq:vam}); under the fair hazard-radius match it reaches 14/15 GO --- one Push seed short of VA-ATACOM's 15/15 (Table~\ref{tab:main}), illustrating that the linearisation almost captures the right behaviour but leaves a residual gap that only the exact quadratic constraint closes. We characterise it here for completeness --- its two tunable knobs ($\alpha$ and the BRT horizon $h$) and their interaction with the control rate. These findings motivated, and are superseded by, the parameter-free velocity-augmented manifold of the main paper.

\subsection{Ablation of the two components}

\begin{table}[t]
\centering
\caption{DT-margin ablation on Goal (5 seeds).}
\label{tab:ablation}
\begin{tabular}{lcccc}
\toprule
Config & $\alpha$ & $h$ & Mean $\bar{C}$ & GO/5 \\
\midrule
Velocity-adaptive only & 0.3 & 0 & 2.90 & 5/5 \\
BRT-only & 0.0 & 3 & 2.35 & 4/5 \\
\textbf{DT-margin (full)} & 0.3 & 3 & \textbf{0.96} & \textbf{5/5} \\
\bottomrule
\end{tabular}
\end{table}

% =====================================================================
% Figure 3 — §V-C ablation bar chart.
%
% Drop-in include for the AAAI 2027 paper main body.  Generated by
% experiments/diagnostics/F3_ablation_bar/plot_f3.py from
% 15 Phase-2 run_eval.json files (3 configs × 5 seeds × 200K env steps).
% Native canvas is 5.2 × 3.0 in.
%
% Usage in §V-C:
%   % =====================================================================
% Figure 3 — §V-C ablation bar chart.
%
% Drop-in include for the AAAI 2027 paper main body.  Generated by
% experiments/diagnostics/F3_ablation_bar/plot_f3.py from
% 15 Phase-2 run_eval.json files (3 configs × 5 seeds × 200K env steps).
% Native canvas is 5.2 × 3.0 in.
%
% Usage in §V-C:
%   % =====================================================================
% Figure 3 — §V-C ablation bar chart.
%
% Drop-in include for the AAAI 2027 paper main body.  Generated by
% experiments/diagnostics/F3_ablation_bar/plot_f3.py from
% 15 Phase-2 run_eval.json files (3 configs × 5 seeds × 200K env steps).
% Native canvas is 5.2 × 3.0 in.
%
% Usage in §V-C:
%   \input{figures/fig_v_c_ablation_bar/fig_v_c_ablation_bar.tex}
%
% Cross-references: Table~\ref{tab:ablation}, §V-C narrative.
% =====================================================================
\begin{figure}[t]
  \centering
  \includegraphics[width=\columnwidth]%
    {figures/fig_v_c_ablation_bar/fig_v_c_ablation_bar.pdf}
  \caption{%
    \textbf{Ablation: both components are necessary.}
    Mean episode cost $\bar C$ on \texttt{SafetyPointGoal1-v0} for
    three configurations (5 seeds $\times$ 200K env steps each); error
    bars are $\pm$1 standard deviation, dots are per-seed values, dashed
    green line is the GO threshold $\bar C \le 5$.
    \emph{VAM-only} ($\alpha{=}0.3$, $h{=}0$) handles M1 chord excursion
    but lacks lookahead for M2 drift, giving 5/5 GO with consistent but
    elevated cost (2.90).
    \emph{BRT-only} ($\alpha{=}0.0$, $h{=}3$) anticipates M2 drift in
    4 seeds (zero or near-zero cost) but a single seed collapses to
    $\bar C{=}7.98$ when momentum at the keepout boundary triggers
    M1 chord excursion before the BRT lookahead activates.
    \emph{DT-ATACOM} ($\alpha{=}0.3$, $h{=}3$) combines both mechanisms,
    reaching 5/5 GO with the lowest mean (0.96) and lowest variance.%
  }
  \label{fig:v_c_ablation_bar}
\end{figure}

%
% Cross-references: Table~\ref{tab:ablation}, §V-C narrative.
% =====================================================================
\begin{figure}[t]
  \centering
  \includegraphics[width=\columnwidth]%
    {figures/fig_v_c_ablation_bar/fig_v_c_ablation_bar.pdf}
  \caption{%
    \textbf{Ablation: both components are necessary.}
    Mean episode cost $\bar C$ on \texttt{SafetyPointGoal1-v0} for
    three configurations (5 seeds $\times$ 200K env steps each); error
    bars are $\pm$1 standard deviation, dots are per-seed values, dashed
    green line is the GO threshold $\bar C \le 5$.
    \emph{VAM-only} ($\alpha{=}0.3$, $h{=}0$) handles M1 chord excursion
    but lacks lookahead for M2 drift, giving 5/5 GO with consistent but
    elevated cost (2.90).
    \emph{BRT-only} ($\alpha{=}0.0$, $h{=}3$) anticipates M2 drift in
    4 seeds (zero or near-zero cost) but a single seed collapses to
    $\bar C{=}7.98$ when momentum at the keepout boundary triggers
    M1 chord excursion before the BRT lookahead activates.
    \emph{DT-ATACOM} ($\alpha{=}0.3$, $h{=}3$) combines both mechanisms,
    reaching 5/5 GO with the lowest mean (0.96) and lowest variance.%
  }
  \label{fig:v_c_ablation_bar}
\end{figure}

%
% Cross-references: Table~\ref{tab:ablation}, §V-C narrative.
% =====================================================================
\begin{figure}[t]
  \centering
  \includegraphics[width=\columnwidth]%
    {figures/fig_v_c_ablation_bar/fig_v_c_ablation_bar.pdf}
  \caption{%
    \textbf{Ablation: both components are necessary.}
    Mean episode cost $\bar C$ on \texttt{SafetyPointGoal1-v0} for
    three configurations (5 seeds $\times$ 200K env steps each); error
    bars are $\pm$1 standard deviation, dots are per-seed values, dashed
    green line is the GO threshold $\bar C \le 5$.
    \emph{VAM-only} ($\alpha{=}0.3$, $h{=}0$) handles M1 chord excursion
    but lacks lookahead for M2 drift, giving 5/5 GO with consistent but
    elevated cost (2.90).
    \emph{BRT-only} ($\alpha{=}0.0$, $h{=}3$) anticipates M2 drift in
    4 seeds (zero or near-zero cost) but a single seed collapses to
    $\bar C{=}7.98$ when momentum at the keepout boundary triggers
    M1 chord excursion before the BRT lookahead activates.
    \emph{DT-ATACOM} ($\alpha{=}0.3$, $h{=}3$) combines both mechanisms,
    reaching 5/5 GO with the lowest mean (0.96) and lowest variance.%
  }
  \label{fig:v_c_ablation_bar}
\end{figure}


Table~\ref{tab:ablation} (Fig.~\ref{fig:v_c_ablation_bar}) shows both DT-margin components contribute: the velocity-adaptive margin alone reaches 5/5 GO but at higher variance (2.90 vs.\ 0.96); BRT shaping alone gives 4/5 GO (one seed fails at $C=7.98$ via M1); the full heuristic attains the lowest cost (0.96) and 100\% GO.

\subsection{$\Delta t$ sensitivity}
\label{sec:dt-sweep}

\begin{table}[t]
\centering
\caption{DT-margin cost vs.\ time step ($\bar C \pm \sigma$, 5 seeds, 50K steps).}
\label{tab:dt}
\begin{tabular}{lccc}
\toprule
$\Delta t$ & ATACOM & DT-margin ($h{=}3$) & Note \\
\midrule
0.01s  & 0.00  & 3.02 & ATACOM OK \\
0.02s  & 0.00  & 0.00 & ATACOM OK \\
0.05s  & 39.95 & 4.96$\pm$6.08 (3/5 GO) & Basin regime \\
0.10s  & 13.68 & \textbf{2.60} & DT-margin OK \\
0.20s  &  ---  & 1.33$\pm$1.50 & $h\!\gg\!h^\star$ \\
\bottomrule
\end{tabular}
\end{table}

% =====================================================================
% Figure 1 — §IV-A Δt phase transition.
%
% Drop-in include for the AAAI 2027 paper main body.  Generated by
% experiments/diagnostics/D5_dt_sweep/plot_d5.py from 60
% Phase-3 D5 result.json files (5 filters × 4 Δt × 3 seeds × 50K env
% steps).  Native canvas is 5.2 × 3.6 in.
%
% Usage in the body of §IV-A:
%
%   % =====================================================================
% Figure 1 — §IV-A Δt phase transition.
%
% Drop-in include for the AAAI 2027 paper main body.  Generated by
% experiments/diagnostics/D5_dt_sweep/plot_d5.py from 60
% Phase-3 D5 result.json files (5 filters × 4 Δt × 3 seeds × 50K env
% steps).  Native canvas is 5.2 × 3.6 in.
%
% Usage in the body of §IV-A:
%
%   % =====================================================================
% Figure 1 — §IV-A Δt phase transition.
%
% Drop-in include for the AAAI 2027 paper main body.  Generated by
% experiments/diagnostics/D5_dt_sweep/plot_d5.py from 60
% Phase-3 D5 result.json files (5 filters × 4 Δt × 3 seeds × 50K env
% steps).  Native canvas is 5.2 × 3.6 in.
%
% Usage in the body of §IV-A:
%
%   \input{figures/fig_iv_a_dt_sweep/fig_iv_a_dt_sweep.tex}
%
% Width selection:
%   - default below is `figure` (single column, scales to \columnwidth)
%   - swap to `figure*` for full text width if needed
%
% Cross-references: Prop. 1 (M1) and §IV-A's empirical claim.
% =====================================================================
\begin{figure}[t]
  \centering
  \includegraphics[width=\columnwidth]%
    {figures/fig_iv_a_dt_sweep/fig_iv_a_dt_sweep.pdf}
  \caption{%
    \textbf{Filter-zoo phase transition at $\Delta t\!\approx\!0.05$~s.}
    Mean episode cost $\bar C$ versus control step $\Delta t$ on
    \texttt{SafetyPointGoal1-v0} (5 filters $\times$ 4 $\Delta t$
    values $\times$ 3 seeds $\times$ 50K env steps).  Vertical bars
    show $\pm$1 standard deviation across seeds.  The green band is
    the GO threshold ($\bar C \le 5$, the convention used throughout
    §III-A).
    \emph{Red, ATACOM}: zero cost at $\Delta t \le 0.02$~s (the
    regime where the continuous-time forward-invariance theorem of
    §III-B holds empirically) and catastrophic at
    $\Delta t \ge 0.05$~s — a clean phase transition consistent
    with the M1 threshold (Prop.~1, §IV-A).
    \emph{Purple, ATACOM-LA}: uniformly poor across $\Delta t$ — the
    1-step lookahead does not save the family and introduces over-
    conservatism artefacts (§VI-A).
    \emph{Orange, DistAdapt} (no lookahead) and \emph{grey, PPO
    baseline}: the velocity-margin alone is insufficient; the
    no-filter baseline is uniformly catastrophic.
    \emph{Blue, Ours} (DistAdapt + sim-BRT): stays in the GO band
    at every $\Delta t$ except the corner cell at $\Delta t=0.05$~s
    where $h\!\cdot\!\Delta t=0.15$~s straddles the M2 drift
    timescale (discussed in §VI-C).%
  }
  \label{fig:iv_a_dt_sweep}
\end{figure}

%
% Width selection:
%   - default below is `figure` (single column, scales to \columnwidth)
%   - swap to `figure*` for full text width if needed
%
% Cross-references: Prop. 1 (M1) and §IV-A's empirical claim.
% =====================================================================
\begin{figure}[t]
  \centering
  \includegraphics[width=\columnwidth]%
    {figures/fig_iv_a_dt_sweep/fig_iv_a_dt_sweep.pdf}
  \caption{%
    \textbf{Filter-zoo phase transition at $\Delta t\!\approx\!0.05$~s.}
    Mean episode cost $\bar C$ versus control step $\Delta t$ on
    \texttt{SafetyPointGoal1-v0} (5 filters $\times$ 4 $\Delta t$
    values $\times$ 3 seeds $\times$ 50K env steps).  Vertical bars
    show $\pm$1 standard deviation across seeds.  The green band is
    the GO threshold ($\bar C \le 5$, the convention used throughout
    §III-A).
    \emph{Red, ATACOM}: zero cost at $\Delta t \le 0.02$~s (the
    regime where the continuous-time forward-invariance theorem of
    §III-B holds empirically) and catastrophic at
    $\Delta t \ge 0.05$~s — a clean phase transition consistent
    with the M1 threshold (Prop.~1, §IV-A).
    \emph{Purple, ATACOM-LA}: uniformly poor across $\Delta t$ — the
    1-step lookahead does not save the family and introduces over-
    conservatism artefacts (§VI-A).
    \emph{Orange, DistAdapt} (no lookahead) and \emph{grey, PPO
    baseline}: the velocity-margin alone is insufficient; the
    no-filter baseline is uniformly catastrophic.
    \emph{Blue, Ours} (DistAdapt + sim-BRT): stays in the GO band
    at every $\Delta t$ except the corner cell at $\Delta t=0.05$~s
    where $h\!\cdot\!\Delta t=0.15$~s straddles the M2 drift
    timescale (discussed in §VI-C).%
  }
  \label{fig:iv_a_dt_sweep}
\end{figure}

%
% Width selection:
%   - default below is `figure` (single column, scales to \columnwidth)
%   - swap to `figure*` for full text width if needed
%
% Cross-references: Prop. 1 (M1) and §IV-A's empirical claim.
% =====================================================================
\begin{figure}[t]
  \centering
  \includegraphics[width=\columnwidth]%
    {figures/fig_iv_a_dt_sweep/fig_iv_a_dt_sweep.pdf}
  \caption{%
    \textbf{Filter-zoo phase transition at $\Delta t\!\approx\!0.05$~s.}
    Mean episode cost $\bar C$ versus control step $\Delta t$ on
    \texttt{SafetyPointGoal1-v0} (5 filters $\times$ 4 $\Delta t$
    values $\times$ 3 seeds $\times$ 50K env steps).  Vertical bars
    show $\pm$1 standard deviation across seeds.  The green band is
    the GO threshold ($\bar C \le 5$, the convention used throughout
    §III-A).
    \emph{Red, ATACOM}: zero cost at $\Delta t \le 0.02$~s (the
    regime where the continuous-time forward-invariance theorem of
    §III-B holds empirically) and catastrophic at
    $\Delta t \ge 0.05$~s — a clean phase transition consistent
    with the M1 threshold (Prop.~1, §IV-A).
    \emph{Purple, ATACOM-LA}: uniformly poor across $\Delta t$ — the
    1-step lookahead does not save the family and introduces over-
    conservatism artefacts (§VI-A).
    \emph{Orange, DistAdapt} (no lookahead) and \emph{grey, PPO
    baseline}: the velocity-margin alone is insufficient; the
    no-filter baseline is uniformly catastrophic.
    \emph{Blue, Ours} (DistAdapt + sim-BRT): stays in the GO band
    at every $\Delta t$ except the corner cell at $\Delta t=0.05$~s
    where $h\!\cdot\!\Delta t=0.15$~s straddles the M2 drift
    timescale (discussed in §VI-C).%
  }
  \label{fig:iv_a_dt_sweep}
\end{figure}


Table~\ref{tab:dt} and Fig.~\ref{fig:iv_a_dt_sweep} corroborate the premise of \S\ref{sec:failure} (and hence Prop.~4): continuous-time ATACOM achieves zero cost at $\Delta t \le 0.02$s but fails at $\Delta t \ge 0.05$s, the phase transition matching the M1 threshold $\Delta t\|u\| \gtrsim \sqrt{d_\text{safe}\cdot r}$. The DT-margin heuristic tracks the safe regime at $\Delta t \in \{0.10, 0.20\}$s with $h{=}3$. At $\Delta t{=}0.05$s the picture is nuanced: the horizon rule (\ref{eq:h_cond}) prescribes $h^\star{=}4$, but a 5-seed sweep finds $h{=}3$ at 3/5 GO (bimodal), $h{=}4$ at 0/5, and $h{=}5$ at 3/5 (one seed at 37.05). We attribute this to a \emph{BRT-shaping basin}: near the penetration horizon the lookahead reports unavoidable penetration along multiple directions every step, integrating spurious penalty into the policy gradient and trapping training.

\subsection{$\alpha$/$h$ sensitivity}
\label{sec:alpha-h-sweep}

Sweeping $\alpha \in \{0, 0.05, 0.1, 0.3, 0.5\}$ and $h \in \{1, 2, 3, 5\}$ on Goal (5 seeds $\times$ 200K):

\begin{table}[t]
\centering
\caption{DT-margin $\alpha$/$h$ sensitivity on Goal (5 seeds $\times$ 200K).}
\label{tab:alpha-h}
\begin{tabular}{lcccc}
\toprule
Config & $\alpha$ & $h$ & mean $\bar C$ & GO/5 \\
\midrule
no margin            & 0.0  & 2 & 3.37$\pm$2.81 & 4/5 \\
short horizon        & 0.05 & 1 & 1.94$\pm$1.78 & 4/5 \\
both at minimum      & 0.05 & 2 & 2.36$\pm$2.60 & 4/5 \\
above minimum        & 0.1  & 2 & \textbf{1.30$\pm$1.66} & \textbf{5/5} \\
$\alpha$ generous    & 0.3  & 1 & 1.72$\pm$1.46 & 5/5 \\
\textbf{paper choice}& 0.3  & 3 & \textbf{0.96$\pm$0.68} & \textbf{5/5} \\
over-provisioned     & 0.5  & 5 & 3.00$\pm$3.38 & 4/5 \\
\bottomrule
\end{tabular}
\end{table}

Only configurations at least one step above the minimum on \emph{both} axes ($\alpha{\ge}0.1$, $h{\ge}2$) reach 5/5 GO; setting $\alpha{=}0$ drops to 4/5 with a seed at $\bar C{=}8.46$. The over-provisioned config ($\alpha{=}0.5$, $h{=}5$) \emph{regresses} to 4/5 ($\bar C{=}9.26$ outlier) --- large-$h$ BRT shaping can integrate spurious penalty (\S\ref{sec:dt-sweep}) and destabilise training. The DT-margin recipe is thus robust only with margin on both knobs; VA-ATACOM eliminates $\alpha$ entirely, replacing it with the calibrated $a_{\max}$.
